problem:  
Let \( (M^4, g_0) \) be a smooth, closed, oriented 4‑manifold with positive isotropic curvature.  
Assume that there exists a sequence of **Ricci flows with surgery** \( (M, g_k(t))_{t>0} \), with surgery parameters \(\delta_k \to 0\), which are non‑collapsing for all \(t>0\) and converge (in the sense of Bamler–Kleiner singular Ricci flows) to a limiting singular Ricci flow \((\mathcal{X}, g_\infty(t))\).  
Let \(\mathcal{X}_{\mathrm{reg}}\) denote the regular set and \(\mathcal{S} = \mathcal{X} \setminus \mathcal{X}_{\mathrm{reg}}\) the singular set at a fixed time \(t_*>0\).

For each \(k\), define the **L² curvature energy measure**
\[
\eta_k := |{\rm Rm}(g_k(t_*))|^2 \, dV_{g_k(t_*)}.
\]
Suppose (possibly after passing to a subsequence) that
\[
\eta_k \rightharpoonup \eta_\infty
\]
as Radon measures on \(\mathcal{X}\).  

**Question.** —  
Does the limit measure \(\eta_\infty\) admit a canonical decomposition of the form  
\[
\eta_\infty = |{\rm Rm}(g_\infty(t_*))|^2\, dV_{g_\infty(t_*)} + \nu_{\mathrm{sing}},
\]
where the **defect measure** \( \nu_{\mathrm{sing}} \ge 0 \) is supported on the singular set \(\mathcal{S}\)?  

If such a decomposition exists, what is the “dimension’’ of curvature concentration:  
is \(\nu_{\mathrm{sing}}\) absolutely continuous with respect to the 2‑dimensional Hausdorff measure \(\mathcal{H}^2|_{\mathcal{S}}\), suggesting a codimension‑two energy concentration analogous to bubbling in harmonic map flow, or can curvature concentration in 4D Ricci flow occur on higher‑dimensional sets?  

Why it is a "good" problem:  
This refined version isolates the **fundamental question of energy concentration** in four‑dimensional Ricci flow—whether curvature blow‑up behaves like codimension‑two “bubbling’’ phenomena familiar from critical geometric PDEs, or whether a genuinely new type of higher‑dimensional singular set can form.  

Unlike the scalar curvature or entropy formulations, this problem uses the \(L^2\)-curvature measure \(|{\rm Rm}|^2 dV\), which is the natural energy for 4‑manifolds and whose behavior governs compactness theorems, removable‑singularity phenomena, and analytic regularity in many geometric flows.  

A decomposition
\[
\eta_\infty = |{\rm Rm}|^2 + \nu_{\mathrm{sing}}
\]
would directly link weak Ricci flows to techniques developed in Yang–Mills theory and harmonic map flow, where defect measures capture energy loss through bubbling. Determining the Hausdorff dimension of \(\nu_{\mathrm{sing}}\) would identify the *true analytic codimension* of singularities in 4D Ricci flow—a problem bridging Riemannian geometry, geometric measure theory, and nonlinear PDE.  

Thus, the problem is simple to state (“what is the defect measure of curvature concentration?”), yet its resolution would mark a conceptual unification of Ricci flow singularities with the broader theory of analytic curvature defect and bubbling—offering a narrow entrance but profound analytic and geometric depth.